\documentclass[a4paper]{article}
\input{./preamble.tex}

\begin{document}

\problemset
    {Ma~100~~Introduction~to~Math}
    {Spring 2024}
    {Sample \LaTeX\ Showcase}
    {Eric Lee}
    {2024-01-01}

\section{Sample Showcase}

\begin{note}[Sample Comment]
    All of the environments are typically used in problem sets or notes.

    Corrections:
    $$
        \correct{x = 2}{x = 5}
    $$
\end{note}

% PROBLEM 1
\begin{problem}[Introduction]
    Solve for $x$ in each of the following:
    \begin{enumerate}[label=\alph*)]
        \item $x = 17 + \frac{1}{2}$
        \item $x + 1 = x^2 + x + 1$
    \end{enumerate}
\end{problem}

\begin{lineproof}[1a]
    $$
        x = 17 + \frac{1}{2} = \frac{35}{2}
    $$
\end{lineproof}

\begin{lineproof}[1b]
    \begin{align*}
        \cancel{x + 1} &= x^2 + \cancel{x + 1} \\
        \implies 0 &= x^2 \\
        \implies x &= 0
    \end{align*}
\end{lineproof}

\begin{eg}[Problem]
    Details
\end{eg}
\begin{boxexplanation}
    Answer
\end{boxexplanation}

\begin{lemma}[Name]
    Problem
\end{lemma}
\begin{boxproof}
    Many of these types of paired environments are listed in the preamble, and can be used separately if desired.
\end{boxproof}

\begin{remark}
    Single line theorem environments are available
\end{remark}

\begin{code}[Sample Code Formatting]\,
    \begin{lstlisting}[language=C]

        #include <stdio.h>

        int main() {
            printf("Hello World!\n");
            return 0;
        }
    \end{lstlisting}
\end{code}

\end{document}
